\clearpage
\VerdierPhysicalPageAnchor{0215}%
\begingroup\small\setlength{\parskip}{0pt}\setlength{\emergencystretch}{1.5em}
\setlength{\abovedisplayskip}{4pt plus 2pt minus 2pt}%
\setlength{\belowdisplayskip}{4pt plus 2pt minus 2pt}%

\VerdierSemanticAnchor{ch3d-3-2-7-cont}
さらに、(1.2.8) で圏 $D(\mathcal A)^o$ を圏
$D(\mathcal A^o)$ と同一視した。これにより、集合
$\operatorname{Ext}_{\mathcal A^o}^{n}(Y,X)$ と
$\operatorname{Ext}_{\mathcal A}^{n}(X,Y)$ を同一視できる
（第 I 章，1.3.9）。これらの同一視のもとで、始域と終域がともに同じ
二つの写像を得る：

\VerdierSemanticAnchor{ch3d-3-2-7-2}
\begin{equation}
\resizebox{0.74\textwidth}{!}{%
$\begin{tikzcd}[ampersand replacement=\&,column sep=5.0em,row sep=2.8em]
E_{\mathcal A}^{n}(X,Y)
  \arrow[r,"{\delta_{\mathcal A}^{n}}"]
  \arrow[d,equal]
\& \operatorname{Ext}_{\mathcal A}^{n}(X,Y)
  \arrow[d,equal] \\
E_{\mathcal A^o}^{n}(Y,X)
  \arrow[r,"{\delta_{\mathcal A^o}^{n}}"']
\& \operatorname{Ext}_{\mathcal A^o}^{n}(Y,X)
\end{tikzcd}$%
}
\tag{3.2.7.2}
\end{equation}
そして次の関係がある：

\VerdierSemanticAnchor{ch3d-3-2-7-3}
\begin{equation}
\delta_{\mathcal A}^{n}=\delta_{\mathcal A^o}^{n}.
\tag{3.2.7.3}
\end{equation}

\VerdierSemanticAnchor{ch3d-3-2-7-proof}
この関係は $n=1$ のとき成り立つ (1.3.4)。そこで (3.2.5) を用いて
$n$ に関する帰納法を行えば証明される。

\VerdierSemanticAnchor{ch3d-3-2-8}
\medskip\noindent\textbf{3.2.8.}\quad
次を
\[
S^n(X,Y)=0\longrightarrow Y\xrightarrow{u_n}Z_{n-1}
\xrightarrow{u_{n-1}}\cdots\xrightarrow{u_2}Z_1
\xrightarrow{u_1}Z_0\xrightarrow{u_0}X\longrightarrow0
\]
$E_{\mathcal A}^{n}(X,Y)$ の元とし、
$f\colon X'\longrightarrow X$ と $g\colon Y\longrightarrow Y''$ を
$\mathcal A$ の二つの射とする。次のようにおく：

\VerdierSemanticAnchor{ch3d-3-2-8-1}
\begin{equation}
\begin{aligned}
S^n(X,Y)*f={}&0\longrightarrow Y\xrightarrow{u_n}Z_{n-1}
\xrightarrow{u_{n-1}}\cdots\xrightarrow{u_2}Z_1
\xrightarrow{u'_1}Z'_0\\
&\xrightarrow{u'_0}X'\longrightarrow0,
\end{aligned}
\tag{3.2.8.1}
\end{equation}
ここで $Z'_0$ はファイバー積である：

\VerdierSemanticAnchor{ch3d-3-2-8-pullback-square}
\begin{center}
\begin{tikzcd}[column sep=4.2em,row sep=2.8em]
Z'_0 \arrow[r,"u'_0"] \arrow[d,"f'"']
  & X' \arrow[d,"f"] \\
Z_0 \arrow[r,"u_0"'] & X
\end{tikzcd}
\end{center}
また $u'_1$ は等式 $f'u'_1=u_1$ と $u'_0u'_1=0$ によって一意に
定まる射である。同様に次のようにおく：

\VerdierSemanticAnchor{ch3d-3-2-8-2}
\begin{equation}
\begin{aligned}
g*S^n(X,Y)={}&0\longrightarrow Y''\xrightarrow{u''_n}Z''_{n-1}
\xrightarrow{u''_{n-1}}Z_{n-2}\xrightarrow{u_{n-2}}\cdots\\
&\cdots\xrightarrow{u_1}Z_0\xrightarrow{u_0}X\longrightarrow0.
\end{aligned}
\tag{3.2.8.2}
\end{equation}
\endgroup
